\section{Quadratic signatures and balanced summation}
\label{sec:signatures}

The transport construction uses quadratic forms with at most one positive
eigenvalue. We prove the required property for integer coordinates and then
show that it survives summation over balanced slots. Integer coordinates
allow two units to be removed at the same location, so the diagonal entries
and the endpoints of the summation require explicit treatment. Large
capacities occur only in finite mathematical sums in this section; the
algorithm will not enumerate those sums.
The one-positive-eigenvalue viewpoint is central to the theory of
Lorentzian polynomials developed by Br\"and\'en and Huh~\cite{BH2020}.
Here the needed closure is a discrete balanced-pair summation, established
by the laminar decomposition and telescoping calculation below.

\begin{definition}[Softened weights]\label{def:soft-weight}
Let $\mathcal B$ and $\mathcal S$ be the large and small slots from
Section~\ref{sec:graph}, and put
\[
 T=\sum_s B_s,\qquad
 \mathcal Q=\prod_s\{0,\ldots,B_s\}^2.
\]
For a full coordinate vector $\xi=(x_s,y_s)_s$, write
$|\xi|=\sum_s(x_s+y_s)$ and $q_s=B_s-y_s$.
For $a\in\R$, let $a_+=\max(a,0)$.  Define
\begin{align*}
 D(\xi)={}&\sum_i\left|\sum_jx_{ij}-\bar r_i\right|
       +\sum_j\left|\sum_iq_{ij}-\bar c_j\right|\\
 &+\sum_i\left(\sum_{j:(i,j)\in\mathcal S}x_{ij}-r_i\right)_+
       +\sum_j\left(\sum_{i:(i,j)\in\mathcal S}q_{ij}-c_j\right)_+.
\end{align*}
For $0<\eta<1$, set
\[
 F_\eta(\xi)=
 \begin{cases}
 \eta^{D(\xi)},&\xi\in\mathcal Q,\quad |\xi|=T,\\
 0,&\text{otherwise}.
 \end{cases}
\]
For a small-coordinate configuration $X$ of total $\sum_{s\in\mathcal S}B_s$,
define
\begin{equation}\label{eq:soft-marginal}
 f_\eta(X)=
 \sum_{\substack{(x_s,y_s)=(a_s,B_s-a_s),\ 0\le a_s\le B_s\\
                  s\in\mathcal B}}
 F_\eta(X,(x_s,y_s)_{s\in\mathcal B}).
\end{equation}
An empty product of choices contributes one term.
\end{definition}

Every box-valid full vector of total $T$ has strictly positive weight,
regardless of its margin deviations or slot occupancies.  Thus
$f_\eta(X)>0$ for every box-valid $X$ of the indicated total.  On the
transversal and defect patterns,
\begin{equation}\label{eq:soft-limit}
 \lim_{\eta\downarrow0}f_\eta(X)=f(X).
\end{equation}
Indeed, a term has $D=0$ exactly when the full row and column margins are
$\bar r,\bar c$ and the small row and column totals do not exceed $r,c$.
When the large slots are balanced, these conditions specify precisely the
nonnegative residuals and padded completions defining $f$.  On a line
without a large slot, the full-margin equality also makes the residual
zero.  The capacity $C$ excludes no such completion, by its choice in
\eqref{eq:parameters}.  Each zero-penalty term tends to one, all other
terms tend to zero, and the sum is finite.

We use the following convention throughout.  Some coordinates have fixed
box-valid values $\sigma$, of total $c=|\sigma|$; a disjoint collection
$\mathcal R$ of whole slots is summed over its balanced profiles; and all
remaining coordinates are active.  For an assignment $v$ to the active
coordinates, let
\[
 \Phi_{\sigma,\mathcal R}(v)
   =\sum_{b\in\prod_{s\in\mathcal R}
                  \{(a,B_s-a):0\le a\le B_s\}}
        F_\eta(\sigma,v,b).
\]
Thus the summation domain is fixed, even when $v$ changes.
A box-valid active display $p$ for two removals satisfies
\begin{equation}\label{eq:two-removal-total}
 c+|p|+\sum_{s\in\mathcal R}B_s=T+2.
\end{equation}
Its two-removal matrix has entries
$M_{ih}=\Phi_{\sigma,\mathcal R}(p-e_i-e_h)$.
In particular, two removals at the same coordinate give a diagonal entry;
an invalid removal gives zero.  The display has total $T+2$, whereas
every nonzero entry is evaluated at total $T$.

\begin{lemma}[Laminar factors]\label{lem:laminar-signature}
With no balanced slots summed, every two-removal matrix just defined has
at most one positive eigenvalue.
\end{lemma}
\begin{proof}
Before imposing the box and total, the formula $\eta^{D(\xi)}$ is a
product of positive log-concave sequences in sums over a laminar family
of coordinate sets.  Here laminar means that any two sets are disjoint
or one contains the other.  The sets are the $x$ coordinates of each row
and their small-slot subsets, and the $y$ coordinates of each column and
their small-slot subsets.  The row and column families are disjoint
because they use different coordinates.  For example, the column
factors have the form
\[
 \eta^{|K-z-\bar c_j|},\qquad
 \eta^{(K'-z-c_j)_+},
\]
where $z$ is the relevant sum of $y$ coordinates and $K,K'$ are the
corresponding sums of capacities.  All exponents are convex in $z$;
since $\log\eta<0$, the factors are log-concave on all integers.

An active coordinate with $p_i=0$ gives a zero row and column of $M$;
discard these indices temporarily.  Restrict the laminar sets to the
remaining indices and absorb fixed coordinates into shifts of their
scalar factors.  Combine factors whose restricted sets coincide, and
absorb empty sets into a positive constant.  We obtain distinct
nonempty laminar sets $G$, positive log-concave sequences $\psi_G$, and
arguments $t_G$ at the display.  Evaluate their formula even for invalid
removals for the moment.  Let $W>0$ be its value at the display and put
\[
 \alpha_i=\prod_{G\ni i}\frac{\psi_G(t_G-1)}{\psi_G(t_G)},
 \qquad
 \gamma_G=\frac{\psi_G(t_G)\psi_G(t_G-2)}{\psi_G(t_G-1)^2}\in(0,1].
\]
The extended two-removal matrix is
$W\diag(\alpha)K\diag(\alpha)$, where
\[
 K_{ih}=\prod_{G\supseteq\{i,h\}}\gamma_G.
\]
For each set define
\[
 a_G=(1-\gamma_G)\prod_{G'\supsetneq G}\gamma_{G'}\ge0.
\]
The sets containing any given pair of indices form a chain.  Telescoping
the product along this chain gives the exact identity
\begin{equation}\label{eq:laminar-blocks}
 K=\one\one^{\mathsf T}
       -\sum_G a_G\one_G\one_G^{\mathsf T}.
\end{equation}
Consequently $K$ is nonpositive on $\{v:\one^{\mathsf T}v=0\}$ and
has at most one positive eigenvalue.  The only still-invalid removals
are diagonal ones at coordinates with $p_i=1$.  Replacing each such
entry by zero subtracts a nonnegative diagonal matrix from $K$, so preserves
that conclusion.  Positive diagonal congruence and restoration of the
zero rows and columns complete the proof.
\end{proof}

We next sum over a fresh slot of capacity $B$.  Two occupancies are
needed.  At occupancy $B$, summing the profiles $(j,B-j)$ eliminates a
balanced slot and preserves the signature.  At occupancy $B+1$, the
same calculation gives a weighted quadratic estimate; this is the
estimate used to bound the transport construction in the next section.
Both conclusions follow from a discrete second-difference inequality.

\begin{lemma}[Pair summation]\label{lem:pair-summation}
Fix a box-valid display $p$ on the old active coordinates, a fresh slot
of capacity $B\ge2$, and fixed context and balanced sums as above.
Choose its displayed occupancy $D'\in\{B,B+1\}$ so that
\begin{equation}\label{eq:pair-total}
 c+|p|+D'+\sum_{s\in\mathcal R}B_s=T+2.
\end{equation}
Write $\Phi(r,a,b)$ for the weight after the balanced sums, with old
values $r$ and fresh pair $(a,b)$.  The box-valid displays
$(p,j,D'-j)$ are indexed by
\[
 \mathcal J=
 \begin{cases}
 \{0,\ldots,D'\},&D'=B,\\
 \{1,\ldots,D'-1\},&D'=B+1.
 \end{cases}
\]
Their two-removal matrices have the blocks
\begin{align*}
 (V_j)_{ih}&=\Phi(p-e_i-e_h,j,D'-j),\\
 S(l)_i&=\Phi(p-e_i,l,D'-1-l),\\
 u_j&=\Phi(p,j-1,D'-1-j).
\end{align*}
Thus $V_j$ records two old removals, $S(l)$ one old and one fresh
removal, and $u_j$ one removal from each fresh coordinate of display
$j$.  Put $S(l)=0$ outside
$0\le l\le D'-1$ and $u_j=0$ outside $1\le j\le D'-1$.
Suppose each matrix
\begin{equation}\label{eq:pair-block}
 M_j=\begin{pmatrix}
 V_j&S(j-1)&S(j)\\
 S(j-1)^{\mathsf T}&u_{j-1}&u_j\\
 S(j)^{\mathsf T}&u_j&u_{j+1}
 \end{pmatrix},\qquad j\in\mathcal J,
\end{equation}
has at most one positive eigenvalue.  The entries $u_{j-1}$ and
$u_{j+1}$ are the double removals at the first and second fresh
coordinates, respectively.
For every real old vector $h$ satisfying
$\sum_{l=0}^{D'-1}h^{\mathsf T}S(l)=0$, define
\begin{equation}\label{eq:pair-cumulative}
 g_j=\frac{\sum_{l=0}^{j-1}h^{\mathsf T}S(l)}{u_j}
       \quad(1\le j\le D'-1),\qquad
 g_j=0\quad\text{otherwise}.
\end{equation}
All these divisions are valid, and
\begin{equation}\label{eq:pair-one}
 -h^{\mathsf T}V_jh\ \ge\
 2u_jg_j^2-u_{j-1}g_{j-1}^2-u_{j+1}g_{j+1}^2,
 \qquad j\in\mathcal J.
\end{equation}
If $D'=B$, the matrix $\sum_{j=0}^B V_j$ has at most one positive
eigenvalue.
If $D'=B+1$, then
\begin{equation}\label{eq:pair-weighted}
 2\sum_{j=1}^{D'-1}u_jg_j^2
 \ \le\
 -h^{\mathsf T}\left(\sum_{j=1}^{D'-1}j(D'-j)V_j\right)h.
\end{equation}
\end{lemma}
\begin{proof}
First, $u_j>0$ for every $1\le j\le D'-1$.  Its pair values
$(j-1,D'-1-j)$ are within $[0,B]^2$ in both cases.  They have total
$D'-2$, so \eqref{eq:pair-total} puts the full configuration on the
total-$T$ slice.  Choose any balanced box-valid profile, for example
$(0,B_s)$, at each summed slot.  The unchanged old profile and fixed
context then give a positive term of $F_\eta$.  This proves positivity
even when the context violates the original margins or some old
displayed coordinates are zero.

The cumulative definition and the balance condition imply, including
$l=0,D'-1$,
\begin{equation}\label{eq:pair-difference}
 h^{\mathsf T}S(l)=u_{l+1}g_{l+1}-u_lg_l.
\end{equation}
We use an elementary consequence of the signature assumption.  If a
symmetric matrix $M$ has at most one positive eigenvalue and
$e^{\mathsf T}Me>0$, then
\begin{equation}\label{eq:positive-projection}
 v^{\mathsf T}Mv\le
       \frac{(v^{\mathsf T}Me)^2}{e^{\mathsf T}Me}.
\end{equation}
Indeed, subtract from $v$ its $M$-orthogonal projection onto $e$.
The remaining vector must have nonpositive square, since otherwise it
and $e$ span a positive definite two-dimensional subspace.

For $j\in\mathcal J$, put $v=(h,g_j,-g_j)$, and let $e_a,e_b$ denote
the two fresh coordinate directions.  Equations
\eqref{eq:pair-block} and \eqref{eq:pair-difference} give
\[
 v^{\mathsf T}M_je_a=u_{j-1}(g_j-g_{j-1}),\qquad
 v^{\mathsf T}M_je_b=u_{j+1}(g_{j+1}-g_j).
\]
If $u_{j-1}>0$, apply \eqref{eq:positive-projection} to $e_a$; if
$u_{j+1}>0$, apply it to $e_b$.  Either bound is at most
\[
 u_{j-1}(g_j-g_{j-1})^2+u_{j+1}(g_j-g_{j+1})^2.
\]
If both diagonals vanish, then $u_j>0$ and $e_a+e_b$ has square $2u_j$
and is $M_j$-orthogonal to $v$.  The same upper bound, now zero, follows.
These cases cover all admissible $j$: at $j=0$ or $D'$ the inward
diagonal is positive, and at every other index $u_j>0$.
An exact expansion now yields
\begin{align*}
 &v^{\mathsf T}M_jv
   -u_{j-1}(g_j-g_{j-1})^2-u_{j+1}(g_j-g_{j+1})^2\\
 &\hspace{2em}=h^{\mathsf T}V_jh
                 +2u_jg_j^2-u_{j-1}g_{j-1}^2-u_{j+1}g_{j+1}^2,
\end{align*}
proving \eqref{eq:pair-one}.

For clarity, the smallest capacity is covered without a positive
diagonal at every display.  When $B=D'=2$, the displays $j=0,2$ use
the diagonal $u_1>0$, and $j=1$ uses the positive sum direction with
square $2u_1$.  When $B=2,D'=3$, the only displays are $j=1,2$;
they have, respectively, positive diagonal $u_2$ and $u_1$.
The out-of-box displays $j=0,3$ are never evaluated.

Write $a_j=u_jg_j^2$, zero outside $1\le j\le D'-1$.
For $D'=B$, summing \eqref{eq:pair-one} over $0\le j\le D'$ gives
\[
 -h^{\mathsf T}\Bigl(\sum_{j=0}^{D'}V_j\Bigr)h
 \ge\sum_{j=0}^{D'}(2a_j-a_{j-1}-a_{j+1})=0.
\]
Thus the sum is nonpositive on the kernel of
$h\mapsto h^{\mathsf T}\sum_l S(l)$, a subspace of codimension at most
one.  A positive eigenspace of dimension two would intersect that
kernel nontrivially, so the sum has at most one positive eigenvalue.
For $D'=B+1$, let $w_j=j(D'-j)$, with $w_0=w_{D'}=0$.
Multiplying \eqref{eq:pair-one} by $w_j$ and summing only over the
admissible indices gives
\[
 \sum_{j=1}^{D'-1}w_j(2a_j-a_{j-1}-a_{j+1})
 =\sum_{j=1}^{D'-1}(2w_j-w_{j-1}-w_{j+1})a_j
 =2\sum_{j=1}^{D'-1}a_j.
\]
This proves \eqref{eq:pair-weighted}, using only scalar zero endpoint
weights, never endpoint matrices outside the box.
\end{proof}

\begin{proposition}[Signature after balanced summation]
\label{prop:signature}
Fix $0<\eta<1$.  For arbitrary box-valid fixed context $\sigma$,
arbitrary balanced-slot summation set $\mathcal R$, and arbitrary
box-valid active display $p$ satisfying \eqref{eq:two-removal-total},
the matrix
\[
 \bigl(\Phi_{\sigma,\mathcal R}(p-e_i-e_h)\bigr)_{i,h}
\]
has at most one positive eigenvalue.
\end{proposition}
\begin{proof}
Induct on $|\mathcal R|$.  Lemma~\ref{lem:laminar-signature} proves
the empty case.  For a nonempty set choose $t\in\mathcal R$, regard
its two coordinates as active at $(j,B_t-j)$, and sum only over
$\mathcal R\setminus\{t\}$.  Every resulting matrix $M_j$ has at most
one positive eigenvalue by induction; its total is $T+2$ by
\eqref{eq:two-removal-total}.  Apply Lemma~\ref{lem:pair-summation}
with $D'=B_t$.  The matrix $\sum_{j=0}^{B_t}V_j$ is exactly the required
matrix after also summing slot $t$, and has the asserted signature.
All capacities here are at least two.  Zero active coordinates give
zero rows and columns throughout and need no additional assumption.
\end{proof}

This is a specific closure under balanced-pair summation, proved by
\eqref{eq:pair-one}; arbitrary sums of matrices with this signature
need not preserve it.  In the transport argument, a weighted matrix
$V_t=\sum_j j(D'-j)V_j$ will be tested against real vectors $h$ whose
entries may have either sign.  Equation~\eqref{eq:pair-weighted}
controls these quadratic evaluations under the stated balance condition.
We make no signature assertion for $V_t$ itself.  In particular,
Proposition~\ref{prop:signature} is applied before this weighted sum,
on the entire softened total-$T$ box, even when the resulting
two-removal configurations lie outside the graph patterns.
